\section{女数学家对曲线的贡献}
\markboth{98. 女数学家对曲线的贡献}{98. 女数学家对曲线的贡献}

玛丽·安妮西（1718～1799）是一位意大利的女数学家，她曾经研究过一种曲线，叫做箕舌线，定义如下：

\par\medskip\noindent
\begin{minipage}[c]{0.46\textwidth}
  \centering
  \begin{tikzpicture}[scale=0.98]
    % 箕舌线 y = a^3 / (x^2 + a^2), 取 a = 1.25
    \def\a{1.25}
    % 坐标轴
    \draw[line width=0.6pt, -{Stealth[length=4pt]}] (-1.85,0) -- (1.85,0);
    \draw[line width=0.6pt, -{Stealth[length=4pt]}] (0,-0.25) -- (0,1.65);

    % 直线 y = a
    \draw[line width=0.6pt] (-1.75,\a) -- (1.75,\a);
    % 圆 x^2 + y^2 = a y
    \draw[line width=0.65pt] (0,{\a/2}) circle ({\a/2});
    % 箕舌线
    \draw[line width=0.9pt, curveblue, domain=-1.72:1.72, samples=120, smooth]
      plot (\x, {\a*\a*\a/(\x*\x + \a*\a)});

    % 射线 OE（斜率 k = 1.35）
    \def\k{1.35}
    \pgfmathsetmacro{\xE}{\a/\k}
    \pgfmathsetmacro{\yE}{\a}
    \pgfmathsetmacro{\xD}{\a*\k/(1+\k*\k)}
    \pgfmathsetmacro{\yD}{\a*\k*\k/(1+\k*\k)}
    \coordinate (O) at (0,0);
    \coordinate (C) at (0,\a);
    \coordinate (E) at (\xE,\yE);
    \coordinate (D) at (\xD,\yD);
    \coordinate (B) at (0,\yD);
    \coordinate (M) at (\xE,\yD);

    \draw[line width=0.55pt] (O) -- ($(O)!1.15!(E)$);
    \draw[line width=0.55pt] (B) -- (M);
    \draw[line width=0.55pt] (E) -- (M);

    \fill (D) circle (1.0pt);
    \fill (E) circle (1.0pt);
    \fill (M) circle (1.1pt);

    \node[below left=-2pt] at (O) {\footnotesize $O$};
    \node[below=1pt] at (1.78,0) {\footnotesize $x$};
    \node[left=1pt] at (0,1.58) {\footnotesize $y$};
    \node[above right=-1pt] at (C) {\footnotesize $C$};
    \node[left=1pt] at (B) {\footnotesize $B$};
    \node[above left=0.5pt, fill=white, inner sep=0.3pt] at (D) {\footnotesize $D$};
    \node[above=1pt] at (E) {\footnotesize $E$};
    \node[right=1pt] at (M) {\footnotesize $M$};
  \end{tikzpicture}
  \par\vspace{3pt}{\footnotesize 图\ 98.1}
\end{minipage}%
\hfill
\begin{minipage}[c]{0.50\textwidth}
  \hspace{2em}如图 98.1，从坐标原点 $O$ 任意作一直线，交圆 $x^2 + y^2 = ay$ 于点 $D$，交直线 $y = a$ 于点 $E$。过 $E$ 作直线平行于 $y$ 轴，过 $D$ 作直线平行于 $x$ 轴，所作两直线的交点为 $M$，则 $M$ 的轨迹叫做箕舌线。
\end{minipage}
\par\medskip

设点 $M$ 的坐标为 $(x, y)$，从图 98.1 可得
\[
  x = BM = CE, \qquad y = OB.
\]
从 $\triangle OBD \backsim \triangle OCE$，得到
\[
  \frac{OB}{OC} = \frac{BD}{CE}.
\]
再注意 $OC = a$，$BD^2 = OB \cdot BC = y(a - y)$，就得到
\[
  \frac{y}{a} = \frac{\sqrt{y(a - y)}}{x}.
\]
化简后，得到箕舌线的方程是
\[
  y = \frac{a^3}{x^2 + a^2}.
\]

在当今世界上，可以找到很多杰出的女数学家；但是在历史上，妇女能够从事数学研究的却寥寥无几。最早的一位女数学家海帕西娅（约 370～415 年）竟然惨遭教会指使的暴徒们残酷杀害。正象海帕西娅出身于数学之家一样，玛丽·安妮西的父亲也是一位数学教授，在意大利波诺格拉大学任教。玛丽从小受到父母对她的丰富多彩的教育，但她不幸是二十一个小孩子中最年长的一个。玛丽·安妮西没有结婚，她把大部分时间用来研究数学和照顾弟弟妹妹，在母亲死后还负担着家务。她用二十年时间钻研各种文字写成的关于微积分的著作，写成两大本教材，被译成英文和法文，广泛地用作课本。

应该指出，安妮西对箕舌线作了详细的研究，但她并不是这种曲线的发现者。在她以前，费马（1657）和格兰弟（1703）就已知道这种曲线。
